\label{sec:conclusion}

\textit{Shelby County v.\ Holder} changed the operational environment
for minority voting rights in redistricting more than any other Supreme
Court decision since \textit{Thornburg v.\ Gingles} established the
three-part test in 1986. By invalidating Section~4(b)'s coverage formula
and rendering Section~5 preclearance a dead letter, the Court removed
the most powerful and most efficient VRA enforcement tool: the ex ante
administrative review that compelled covered jurisdictions to justify
their redistricting plans before implementation.

Thirteen years later, the legislative fix---a new coverage formula---
has not been enacted. Former covered states conduct redistricting
without any federal pre-enactment review. Section~2 vote-dilution
litigation under the \textit{Gingles} framework remains the primary
federal remedy, reinforced by \textit{Allen v.\ Milligan}'s 2023
reaffirmation that the framework applies with full force to redistricting
claims. Several states have enacted VRA analogs that provide supplemental
protection, with New York's John R.\ Lewis Voting Rights Act representing
the most comprehensive restoration of preclearance-like review at the
state level.

Algorithmic redistricting with \textsc{bisect} fills part of the gap
that Section~5's demise created. It cannot restore the ex ante legal
requirement that covered states prove non-retrogression; that requires
Congressional action. But it provides the documentation function---a
transparent, reproducible, verifiable neutral plan that any party can
independently replicate and verify---and a partial deterrence function
through the availability of a published neutral baseline. In Section~2
litigation, the \textsc{bisect} reference plan quantifies the gap between
what geography would produce under a neutral process and what the enacted
map actually produces, translating the legal concept of ``vote dilution''
into a specific, measurable deviation from a neutral baseline.

Three practical conclusions follow for redistricting practitioners
working in the post-\textit{Shelby} environment:

\begin{enumerate}
\item \textbf{Run \textsc{bisect} early}. In former covered states with
  significant minority populations, produce and publish a \textsc{bisect}
  reference plan before the legislative redistricting process begins.
  This establishes the neutral baseline and, by making the deviation from
  that baseline visible, creates a deterrent against extreme redistricting.

\item \textbf{Document the VRA mode analysis}. The \textsc{bisect} VRA mode
  (\texttt{--weights-override vra-aligned}) implements the \textit{Gingles}
  three-prong analysis algorithmically. Run the VRA mode analysis for all
  relevant states and document the results: which districts are majority-minority
  under a neutral plan, what the BVAP thresholds are, and how the enacted
  plan compares. This documentation is the algorithmic equivalent of the
  Section~5 preclearance submission.

\item \textbf{Use the audit chain as evidence}. The \textsc{bisect} SHA
  audit chain provides cryptographic proof that the reference plan was
  produced by the neutral algorithm and has not been modified. In litigation,
  this resolves disputes about whether the algorithmic plan was manipulated
  ex post to support a desired outcome. Present the audit chain as part
  of the expert witness report alongside the demographic and geographic
  analysis.
\end{enumerate}

Section~5 cannot be restored by algorithmic means. But the transparency,
reproducibility, and verifiability that algorithmic redistricting provides
can substitute for the documentation function that Section~5 required,
and can give Section~2 plaintiffs a quantitative, evidence-based framework
for measuring the cost of minority vote dilution in the post-\textit{Shelby}
world.

\bigskip
\noindent\textit{The author thanks the redistricting research team for
data access and algorithmic development. All errors are the author's own.}
